% ================================================================
%  ENCUENTRO 3 — Martes 5 de mayo de 2026
% ================================================================
\chapter{Tercer Encuentro}

\begin{destacado}
    \textbf{Fecha:} Martes, 5 de mayo de 2026 \quad
    \textbf{Horario:} 6:20 a.m.\ -- 8:00 a.m. \quad
    \textbf{Asistentes:} ---
\end{destacado}

\medskip

El encuentro del día de hoy no pudo desarrollarse como estaba planeado. Un cambio de
horario en los grupos de los estudiantes integrantes del semillero impidió la reunión
habitual. El espacio fue aprovechado por el director académico para avanzar en la
preparación de material para próximas sesiones.

% ---------------------------------------------------------------
\section{Situación del Encuentro}
% ---------------------------------------------------------------

Debido a una modificación no prevista en los horarios de clase de los estudiantes
asistentes al semillero, no fue posible llevar a cabo el encuentro según el plan
establecido. Esta situación pone de manifiesto la necesidad de establecer mecanismos
de comunicación ágiles y de contar con un calendario fijo y acordado con las
directivas de la institución para garantizar la regularidad de las sesiones.

\begin{nota}{Encuentro no realizado por cambio de horario}{nota:encuentro3-cancelado}
    El tercer encuentro del semillero no pudo desarrollarse debido a un cambio de
    horario en los grupos de los estudiantes. Se sugiere coordinar con anticipación
    cualquier modificación de horario institucional que pueda afectar la participación
    del semillero, y establecer un protocolo de aviso temprano para todos los
    integrantes.
\end{nota}

% ---------------------------------------------------------------
\section{Trabajo Autónomo del Director Académico}
% ---------------------------------------------------------------

En el tiempo destinado al encuentro, el profesor \textbf{Daniel Soto Villada} aprovechó
la sesión para elaborar un conjunto de problemas de astronomía enfocados en el
\textbf{Sistema Solar}, con énfasis en los tópicos de \textbf{distancias} y
\textbf{tamaños reales}. Estos problemas están diseñados para desarrollar habilidades
de cálculo, conversión de unidades y pensamiento proporcional aplicado al estudio del
cosmos.

\begin{idea}{Problemas de Sistema Solar para próximos encuentros}{idea:problemas-ss}
    El conjunto de problemas elaborado en esta sesión abarca dos grandes tópicos:
    \begin{itemize}[leftmargin=1.5em]
        \item \textbf{Distancias:} ejercicios sobre la Unidad Astronómica (AU),
              escalas de distancia dentro del sistema solar, la Heliosfera y la
              Nube de Oort.
        \item \textbf{Tamaños reales:} problemas comparativos de tamaños de
              planetas, satélites y el Sol, utilizando proporciones y modelos a escala.
    \end{itemize}
    Estos problemas se trabajarán en futuros encuentros, \textbf{después de haber
    presentado el contexto teórico y los conceptos requeridos}, de modo que los
    estudiantes puedan abordarlos con las herramientas necesarias.
\end{idea}

\begin{nota}{Criterio pedagógico de los problemas}{nota:criterio-problemas}
    Los problemas han sido diseñados de forma progresiva: parten de comparaciones
    intuitivas y escala cotidiana, y avanzan hacia cálculos que requieren el manejo
    de potencias de diez, la notación científica y el razonamiento proporcional.
    El objetivo es que los estudiantes desarrollen una intuición real sobre las
    dimensiones del Sistema Solar antes de abordar distancias intergalácticas o
    cosmológicas.
\end{nota}

% ---------------------------------------------------------------
\section{Proyección del Próximo Encuentro}
% ---------------------------------------------------------------

El siguiente encuentro retomará la agenda prevista y se orientará hacia la clase
teórica introductoria sobre el Sistema Solar: contexto, escala de tamaños y
distancias, y marco conceptual previo a la resolución de los problemas preparados.

\begin{compromiso}{Compromisos para la siguiente sesión}{comp:clase3-sig}
    \begin{itemize}[leftmargin=1.5em]
        \item Confirmar con anticipación la disponibilidad de todos los integrantes
              para evitar encuentros sin asistencia.
        \item Presentar la clase teórica introductoria al Sistema Solar (posición
              en el universo, escalas de tamaño y distancias).
        \item Distribuir los problemas de astronomía elaborados en esta sesión para
              que los estudiantes comiencen a familiarizarse con el lenguaje y los
              datos del Sistema Solar.
    \end{itemize}
\end{compromiso}
